\section{ĐỀ SỐ \thesection } % ĐỀ 14 Lê quý đôn q3
\setcounter{bt}{0}

%%Câu 1.
\begin{bt} \
	\begin{enumerate}
		\item Điền chữ số vào dấu * để được số $\overline{4 *}$ là số nguyên tố.
		\item Cho tập hợp $A=\{x \in \mathbb{Z} \mid-4<x<1\}$. Hãy viết tập hợp $A$ bằng cách liệt kê các phần tử.
		\item Hoàn thành bảng dưới đây vào giấy làm bài:
		
		\begin{center}
			\begin{tabular}{|l|c|c|}
				\hline
				Số La Mã & XIV & XIX \\
				\hline
				Giá trị tương ứng trong hệ thập phân & 14 &  \\
				\hline
			\end{tabular}
		\end{center}
		
	\end{enumerate}
	
	\loigiai{
		\begin{enumerate}
			\item  Vậy các giá trị của $x$ để $\overline{4*}$ là số nguyên tố là $1$, $3$, và $7$.
	
			\item $A=\{-3;\ -2;\ -1;\ 0\}$
	
			\item	\begin{tabular}{|l|c|c|}
					\hline
					Số La Mã & XIV & XIX \\
					\hline
					Giá trị tương ứng trong hệ thập phân & 14 & 19 \\
					\hline
				\end{tabular}
	
		\end{enumerate}
	}
\end{bt}


%%Câu 2. 
\begin{bt} 
	\
	\begin{enumerate}
		\item Sắp xếp các số nguyên sau theo thứ tự giảm dần: $0 ;\  2024 ;\ -2023 ;\ -3 ;\ 2$.
		\item Hãy vẽ một trục số rồi biểu diễn các số $-4 ;\ 3 ;\ 1 ;\ -2 ;\ 0$ trên trục số đó.
	\end{enumerate}
	
	\loigiai{
		\begin{enumerate}
			\item Sắp xếp các số nguyên theo thứ tự giảm dần: $2024, 2, 0, -3, -2023$.
			
			\item 

				\begin{tikzpicture}
					% Vẽ trục số
					\draw[->] (-5,0) -- (4,0) node[right] {Trục số};
					
					% Vẽ các số trên trục
					\foreach \x in {-4, -3, -2, -1, 0, 1, 2, 3} {
						\draw (\x, 0.1) -- (\x, -0.1) node[below] {\x};
					}
					
					% Đánh dấu các số cần thiết
					\filldraw[red] (-4,0) circle (2pt);
					\filldraw[red] (-2,0) circle (2pt);
					\filldraw[red] (0,0) circle (2pt);
					\filldraw[red] (1,0) circle (2pt);
					\filldraw[red] (3,0) circle (2pt);
				\end{tikzpicture}

		\end{enumerate}
	}
\end{bt}



%%Câu 3.
\begin{bt}Nhóm các bạn lớp 6A cần chia 54 quyển vở và 48 bút chì vào trong các túi quà để mang tặng các bạn ở cô nhi viện sao cho số quyển vở và bút chì ở mỗi túi đều như nhau. Tính số lượng túi quà nhiều nhất mà nhóm bạn có thể chia được.
	
	\loigiai{Số lượng túi quà nhiều nhất mà nhóm bạn có thể chia được là 6 túi.}
\end{bt}


%%Câu 4.
\begin{bt}
	Tại vòng bảng của một cuộc thi có quy định mỗi thí sinh dự thi phải trả lời 20 câu hỏi. Với mỗi câu trả lời đúng được $+20$ điểm, mỗi câu trả lời sai được $-10$ điểm và 0 điểm cho mỗi câu chưa trả lời. Bạn Hùng trả lời được 12 câu đúng, 3 câu sai và các câu còn lại chưa trả lời được. Hỏi bạn Hùng đạt được bao nhiêu điểm?
	
	\loigiai{Hùng đạt được 210 điểm.}
\end{bt}


%%Câu 5.
\begin{bt} Nhà bạn Nam có một khu vườn hình chữ nhật có chiều dài 25 m, chiều rộng 15 m.
	\begin{enumerate}
		\item Tính chu vi của khu vườn này. 
		\item Mẹ bạn Nam dự định trồng hoa Tết trên khu vườn này. Biết chi phí đầu tư cho mỗi mét vuông trồng hoa Tết là $x$ đồng. Hỏi để trồng hoa Tết mẹ bạn  cần đầu tư bao nhiêu tiền?
	\end{enumerate}
	
	\loigiai{a) Chu vi của khu vườn là $80$ mét \hfill b) $375\cdot x$ đồng.
	}
\end{bt}



%%Câu 6.
\begin{bt} Tìm các số nguyên $x$ sao cho 

	\begin{listEX}[3]
		\item $2(x + 3) =- 16$;
		\item $6x - 4 = 2x + 12$;
		\item $3x^2 - 12 = 0 $.
\end{listEX}
	\loigiai{a) $x=-11$ \hfill b) $x=4$ \hfill c) $x=\pm 2$.
	}
	
	%%Câu 7.
	\begin{bt} Biểu đồ cột dưới đây cho biết thông tin về môn thể thao yêu thích của học sinh lớp 6A.
		\immini{\begin{enumerate}
			\item Hãy đọc thông tin về môn thể thao yêu thích của học sinh lớp 6A trong biểu đồ cột ở trên và lập bảng thống kê tương ứng.
			\item Hãy cho biết môn thể thao nào được nhiều học sinh yêu thích nhất và tính tổng số học sinh của lớp 6A.
			\end{enumerate}
		}{
			\begin{tikzpicture}[scale=1,font=\small,line join=round,>=stealth]
				\def\hoanh{9};
				\def\tung{3.5}; \def\sc{7};
				\node[above] at (\hoanh/2,\tung+0.6){\normalsize \bfseries Môn thể thao yêu thích của học sinh lớp 6A};
				\def\mau{cyan};
				\draw (0,0)node[left]{0}--(\hoanh,0);
				\draw[->] (0,0)--(0,\tung+0.5);
				\foreach \y in {5,10,...,25}{\draw (-0.1,\y/\sc)node[left]{$\y$}--++(0:0.2);
					\draw[dashed,gray] (0,\y/\sc)--(\hoanh-0.2,\y/\sc);}
				\foreach \x/\n in{1/20,3-0.2/7,5-0.4/2,7-0.6/6,9-0.8/10}{\draw[line width=8mm,\mau] (\x,0)--++(0,\n/\sc)node[black,shift={(90:0.25)}]{\n};}
				\foreach \x/\p in {1/Bóng chuyền,3-0.2/Bóng rổ,5-0.4/Đá cầu,7-0.6/Bơi lội,9-0.8/Cầu lông}{\node at (\x,-0.4){\footnotesize \p};}        
			\end{tikzpicture}
	}
			
		
		\loigiai{
			\begin{enumerate}
				\item Bảng thống kê:\\
				\begin{tabular}{|c|c|c|c|c|c|}
						\hline
						Môn thể thao & Bóng chuyền  &Bóng rổ& Đá cầu& Bơi lội& Cầu lông\\
						\hline
						Số học sinh & 20 & 7 & 2 & 6&10 \\ 
						\hline
					\end{tabular}
				\item Môn thể thao được nhiều học sinh yêu thích nhất là bóng chuyền, tổng số học sinh của lớp 6A là 45.
			\end{enumerate} }
	\end{bt}
\end{bt}